% Módulo 02g — Curso embutido: cálculo, floats, estatística, Lean, orquestração (R776+)
\section{Derivação variacional (MaxEnt)}
Maximizar $H(w)=-\sum w_i\ln w_i$ sob $\sum w_i=1$ e média fixada, por multiplicadores de Lagrange, produz $w_i\propto e^{\beta f_i}$: a forma exponencial não é arbitrária, é a distribuição de máxima entropia compatível com a média \cite{jaynes1957,shannon1948}. Com $\beta=1$ recupera-se o softmax; $\beta=1/\tau$ introduz a temperatura. Regras de escore próprias caracterizam quando probabilidades relatadas são honestas \cite{gneiting2007}.
\section{Limites de temperatura}
Com maximizador único $i^*$: quando $\tau\to0^+$, cada razão $e^{(f_j-f_{i^*})/\tau}\to0$ e $w_{i^*}\to1$; quando $\tau\to\infty$, todas as razões tendem a 1 e os pesos equalizam em $1/n$. T10 é a versão finita e verificada dessa intuição: esfriar aumenta estritamente o peso do favorito, sem passar ao limite.
\section{Anatomia do ponto flutuante}
IEEE-754 representa $x\approx m\times2^e$ com arredondamento ao mais próximo; operações elementares têm erro relativo limitado pela unidade de arredondamento, mas exponenciação de argumentos grandes estoura para infinito \cite{goldberg1991,muller2018}. Subtrair o máximo antes de exponenciar mantém todo expoente $\leq0$ (T11), logo cada termo fica em $(0,1]$: o estouro torna-se impossível por construção, e a divisão final reescala sem perda de decisão (T4).
\section{Intervalos, Pearson e pareamento}
Para média $\bar{x}$ de $n$ observações com desvio $s$, o IC95 $\bar{x}\pm1{,}96s/\sqrt{n}$ aproxima a incerteza sob normalidade; o relato por intervalos segue a estatística por estimação \cite{cumming2014} e a advertência contra o uso pontual de valores-\textit{p} \cite{wasserstein2016}; reamostragem bootstrap \cite{efron1979} é a alternativa quando a normalidade é suspeita. Pearson mede relação linear; parear estratégias nas mesmas decisões elimina variação de dificuldade. Regras de escore próprias \cite{gneiting2007} lembram que probabilidade bem relatada é hipótese, não default.
\section{Táticas Lean em prosa}
\textit{Positividade} (\texttt{exp\_pos}, divisão de positivos) estabelece sinais. \textit{Reescrita} (\texttt{sum\_div}, \texttt{sum\_mul}, \texttt{exp\_sub}, \texttt{exp\_add}) move termos entre formas equivalentes. \textit{Monotonicidade} (\texttt{mul\_le\_mul\_of\_nonneg\_left}, \texttt{div\_lt\_div\_iff₀}) preserva desigualdades. \textit{Fechamento} (\texttt{div\_self}, \texttt{linarith}, \texttt{field\_simp}, \texttt{ring}) conclui identidades. \textit{Estrutura} (\texttt{take\_append\_drop}, \texttt{pairwise\_append}, \texttt{add\_sum\_erase}, \texttt{exists\_max\_image}, \texttt{sum\_lt\_sum}, \texttt{sum\_eq\_single\_of\_mem}) decompõe objetos finitos. Cada nome foi conferido na Mathlib pinada, nunca na memória.
\section{Orquestração: do contrato ao quadro-negro}
Alocar tarefas entre resolvedores distribuídos remonta ao protocolo de rede de contratos \cite{smith1980}; integrar conhecimento heterogêneo sob incerteza, aos sistemas de quadro-negro \cite{erman1980}. Nosso roteador instancia a linhagem com meios atuais: gates decidem elegibilidade, cabeças pontuam aspectos, softmax calibra, máximo decide --- cada camada auditável separadamente, como documenta a trilha de pinas e hashes.
